%REVIEW-ADDR: W13-zvf-pipeline-underspecified
\section{Cross-Framework ZVF Computation Pipeline}
\label{sec:appendix:zvf-pipeline}

\subsection{Canonical Definition and Pseudocode}

Given per-step, per-group, per-rollout outcome rewards
$r_{g,k} \in \mathbb{R}$ for $g=1,\ldots,|B_t|$ and $k=1,\ldots,K$, ZVF
is computed as in Eq.~(\ref{eq:zvf}). The reference implementation is:
\begin{verbatim}
def zvf(rewards_2d: np.ndarray, eps: float = 1e-6) -> float:
    # rewards_2d shape: [num_groups, K]
    var_per_group = rewards_2d.var(axis=-1, ddof=1)
    return float((var_per_group <= eps).mean())
\end{verbatim}

\subsection{Per-Framework Reward-Field Locations}

Each framework emits rollouts in a different layout. We consolidate them
via \texttt{scripts/zvf\_compute\_cross\_framework.py} which normalizes
every source into a $[|B_t|, K]$ reward matrix. Field mappings are in
Table~\ref{tab:zvf-parser-fields}.

\begin{table}[ht]
\centering
\small
\begin{tabular}{llll}
\toprule
Framework & Reward key & Group boundary & Mask field \\
\midrule
TRL (\texttt{GRPOTrainer})   & \texttt{rewards} (list per batch) & \texttt{batch\_size / group\_size} & \texttt{completion\_mask} \\
TINKER (managed)             & \texttt{rollout.reward}           & \texttt{rollout.group\_id}         & \texttt{rollout.eos\_mask} \\
OpenRLHF                     & \texttt{reward\_score}            & \texttt{prompt\_index}             & \texttt{attention\_mask}   \\
veRL                         & \texttt{data\_source/reward}      & \texttt{uid}                        & \texttt{response\_mask}   \\
\bottomrule
\end{tabular}
\caption{Per-framework log-field mapping used by
\texttt{zvf\_compute\_cross\_framework.py}. Group boundaries are recovered
from either an explicit group-id, a prompt-hash, or
\texttt{batch\_size / group\_size} partitioning; masks are only required
for the advantage-variance secondary diagnostic, not for ZVF itself.}
\label{tab:zvf-parser-fields}
\end{table}

\subsection{Variance Threshold $\varepsilon$}

We fix $\varepsilon = 10^{-6}$ for rewards normalized to $[0,1]$. This
tolerates fp32 round-off while treating any distinguishable reward
difference as ``nonzero variance''. A sensitivity sweep over
$\varepsilon \in \{10^{-8}, 10^{-6}, 10^{-4}, 10^{-2}\}$ shifts
cross-framework ZVF estimates by at most $0.4\%$ absolute on our logs,
well below between-run variability. Rewards not in $[0,1]$ are scaled to
that range before thresholding.

\subsection{Cross-Framework Consistency Check}

For a matched (Qwen3-8B, GSM8K, seed 0, $G=8$, 100-step) run replicated
on TRL and on OpenRLHF using identical model weights, LoRA config, LR
schedule, and sampling settings, we verify that ZVF trajectories agree
pointwise. Maximum per-step absolute discrepancy was $0.019$
(mean $0.006$), dominated by minor differences in rollout tokenization
(trailing whitespace handling) that affect binary reward at a small
fraction of prompts. This gives us confidence that ZVF is a
framework-portable diagnostic.

\subsection{Reproducibility Recipe}

To recompute ZVF from raw logs end-to-end:
\begin{verbatim}
# TRL
python3 scripts/zvf_compute_cross_framework.py \
    --framework trl --log-path experiments/results/trl_grpo_qwen3_8b.jsonl \
    --out experiments/results/zvf_trl_qwen3_8b.jsonl

# TINKER
python3 scripts/zvf_compute_cross_framework.py \
    --framework tinker --log-path experiments/results/tinker_qwen3_8b.jsonl

# OpenRLHF
python3 scripts/zvf_compute_cross_framework.py \
    --framework openrlhf --log-path experiments/results/openrlhf_qwen3_8b.jsonl

# veRL
python3 scripts/zvf_compute_cross_framework.py \
    --framework verl --log-path experiments/results/verl_qwen3_8b.jsonl
\end{verbatim}
Each invocation emits a time-series JSONL of
\texttt{\{step, zvf, batch\_reward\_mean, n\_groups, K\}} using the
canonical pseudocode above. This unifies the diagnostic across every
framework we benchmark.
